\chapter{Discussion and Conclusion}\label{chap:discussion}\label{chap:conclusion}

This thesis investigates to what extent commentator-inspired engagement strategies, implemented in a data-driven pipeline, can improve the engageability of automated esports live text commentary. Through the design, implementation, and evaluation of the \textit{SAGE} pipeline, we tested two such strategies: determining \textit{what} strategic context to add (\textit{Expert Insights}) and \textit{how} to present events coherently (\textit{Event Structuring}). We focused on live text commentary and evaluated it with end users; this modality shaped expectations and engagement patterns.

In this chapter, we interpret the evaluation findings through our three research questions. We explain why the tested strategies improved skill-centric engagement but not spectacle-centric engagement, situate the results within prior Natural Language Generation (NLG) literature, and derive design implications and validity boundaries.

\section{Interpretation of Findings}

\subsection{RQ1: The Effect of Commentator-Inspired Strategies}
\textbf{RQ1: To what extent do commentator-inspired engagement strategies, implemented in a data-driven pipeline, improve the engageability of automated esports live text commentary?}

Our results show that commentator-inspired strategies do improve engageability, but to a modest and domain-specific degree, with a clear ceiling in our setting. Baseline engageability means were low ($M=2.5$--$3.9$, 7-point scale) and combined scores for many dimensions stayed below the midpoint. Even so, the architectural modules produced moderate, statistically significant effects on specific engageability dimensions. The evaluation (Section~\ref{sec:summary-results}) shows a two-sided pattern: the architectural modules improved \textit{skill-centric} engagement but did not improve \textit{spectacle-centric} engagement.

\subsubsection{Skill vs. Spectacle}
Adding data-driven expert insights yielded significant positive effects ($\alpha = 0.05$) on \textit{Competitive Nature} ($\beta=0.199$), \textit{Skill Improvement} ($\beta=0.186$), and \textit{Vicarious Sensation} ($\beta=0.252$). Low $\beta$ values indicate that while these effects are statistically significant, they are modest in magnitude. These dimensions are closely tied to player skill, and align with the \textit{Domain Knowledge} and \textit{Analytical Knowledge} sources identified in Chapter~\ref{chap:prelim}. However, \textit{Skill Improvement} remained comparatively low in absolute terms.

Conversely, we found no significant effects for \textit{Dramatic Nature} ($\beta=0.103$, $p=0.176$) or \textit{Entertaining Nature} ($\beta=0.112$, $p=0.124$). The \textit{spectacle} of esports, its emotional highs and lows and narrative arcs, may be difficult to capture through data-driven text generation alone. The text-only modality likely contributes, as does the lack of \textit{Background Knowledge} and the anonymization of player names, which can otherwise support dramatic framing.

Baseline \textit{Logicality} was well above the midpoint ($\approx 5.6$), as was \textit{Understanding} ($\approx 4.5$). This indicates that what we called \textit{Descriptive Information} (directly observed game state) can provide coherent summaries of a match, and that \textit{Expert Insights} can increase engageability by adding \textit{Analytical} and \textit{Domain Knowledge}. The A-E-P+S framework and commentator-inspired strategies likely also contributed to the overall logical structure.

\subsubsection{What Participants Said}
Qualitative feedback echoed the statistical findings and pointed to a friction inherent to text as an engagement medium. Participants emphasized that the text-only modality limited the ability to convey excitement, and described a tension between shoutcasting (synchronous, social, exciting) and offline reading of textual live updates (asynchronous, efficient, non-linear). This effect may also be influenced by suspicion of AI-generated content, reflected in the empirical data: AI Sentiment ($M=3.03$, SD$=1.18$, where 4 is neutral) and on average 3.3 out of 4 ($SD=0.77$) texts guessed to be AI-generated.

While not consistent across all outcomes, exploratory correlations also suggest that participants who were more positive about the utility of AI in this context tended to report higher \textit{Game Knowledge} enjoyment (Spearman $\rho=0.53$), however, this question was asked after these outcomes were measured.

Participants also reacted negatively to promotional language ("hype") in the baseline texts. Because \textit{Expert Insights} competes for token space within the Large Language Model (LLM)'s output constraints, it may have displaced some of this content, but the absence of significant interaction effects suggests this was not the primary driver of the observed improvements.

\subsubsection{Skill with a Spectacle Ceiling}
The answer to RQ1 is that commentator-inspired strategies can meaningfully improve engageability, provided the engagement goal is framed as \textit{skill-centric} rather than \textit{spectacle-centric}. However, this conclusion must be tempered by the absolute scores: even our best-performing configuration (\textit{Expert Insights}) rarely exceeded a score of 4.0 (``Neither agree nor disagree'') on many engageability metrics. While the relative improvement was statistically significant, absolute scores never exceeded ``Neutral''.

Despite this ceiling, the direction of the signal is robust. If \textit{Expert Insights} can lift engagement even in a dry, text-only medium, its effect may be amplified in richer media, as suggested by participants in the open fields. Their preference for analysis over description is likely modality-agnostic, implying that this architectural choice is valid for future systems (e.g., live commentary or automated, possibly interactive, analysis of matches). Further gains may come from focusing on skill-centric engagement and reducing promotional language and play-by-play details.

AI systems are unlikely in the near term to fully replicate the Socialization Opportunity and Friends Bonding that human shoutcasters provide. Instead, the textual modality can best support data-driven systems that quickly process high-fidelity data and combine it with external knowledge sources to generate insights that human casters may miss.

\subsection{RQ2: Transforming Data into Insight}
\textbf{RQ2: How can high-fidelity data be transformed into engaging natural language descriptions of complex esports events?}

To understand what kind of data-to-text transformation produces engagement, we look at the performance of the \textit{Expert Insights} module, rather than accuracy of the text representing the underlying data. This suggests that engagement in this context is driven by \textit{contextualisation}, not the \textit{description} of state alone: by adding strategic explanations that interpret the raw data, we significantly improved engagement metrics related to skill.

\subsubsection{The ``Why'' Over the ``What''}
The baseline condition (No Modules) essentially simulated a standard log-based summary: it reported \textit{what} happened (e.g., ``A fight occurred''), often enhanced with promotional language. This was heavily criticized in qualitative feedback as ``boring'' and ``repetitive.'' By contrast, the \textit{Expert Insights} module, which barely changed event descriptions but added strategic context, significantly improved scores on skill-centric engagement metrics.

This contrasts with human-generated soccer commentary \cite{lewandowski2012, pavzanin2022}, where an engaging ``promotional language'' style is observed. In our AI-generated format, participants instead preferred short facts with relevant contextual information: they valued information density and insights over ``floweriness.'' Interpreting \textit{why} the descriptive data matters is what drove engagement.

\subsubsection{More Context, Less Clutter}
The \textit{Detail Level} result shows the same paradox. \textit{Expert Insights} had a significant negative effect on the perceived detail rating ($\beta=-0.243$), moving the system closer to the ideal ``Just Right'' category.

Adding more details of a different type (via \textit{Expert Insights}) made the summary ``feel'' less detail-heavy. Since \textit{Expert Insights} did not increase the available number of tokens in the LLM-realiser, this likely reflects a trade-off from battle detail to contextual detail. Conversely, \textit{Event Structuring}, which filtered out less relevant events, did not have a significant effect on detail level ($\beta=0.064$, n.s.).

This suggests that raw data reporting feels ``cluttered'' because it lacks hierarchy: every skirmish can feel equally important. \textit{Expert Insights} provided context so participants could better understand relevance, reducing the perceived ``noise'' of micro-events. For high-fidelity data, this implies that \textit{summarization requires contextualization, not just reduction}: efficiency means value per word, not word count.

\subsubsection{Technical Implications}
From a technical perspective, this validates our architectural decision to separate \textit{Event Detection} from \textit{Insight Generation}. Our modular design allowed us to inject these insights naturally. However, the qualitative feedback regarding narrative gaps suggests that our underlying event-based abstraction may have been too coarse. While we successfully identified strategic reasons, we often missed the operational context (e.g., \textit{how} the army got there). Furthermore, the reliance on rule-based algorithms for insight generation introduced occasional inaccuracies (e.g., in army size detection during continuous unit production), which users noted.

\subsubsection{Context, Not Description}

High-fidelity data can be used for engaging descriptions when the system prioritizes contextualization over pure event description, and structures information to maximize value efficiency rather than raw detail. In this modality, the practical design target is \textit{value per word}: prioritizing interpretation, relevance, and explicit connections over sheer coverage.

\subsection{RQ3: The Limits of Narrative Structure}
\textbf{RQ3: How can important events in an esports match be selected and structured to form a coherent narrative?}
This is the more technical question, captured mostly by the design of our system, and the \textit{Event Structuring} module in particular. Our system selected events and enriched them with descriptive information (such as the resource balance). The \textit{Event Structuring} module then filtered and reordered these events before realization by the LLM.

Our findings for RQ3 complicate the picture for narrative generation. Hypothesis 2 (\textit{Event Structuring}) was unsupported: filtering events into logical sub-sequences (rather than pure chronology) had no significant effect on any engagement metric (e.g., \textit{Understanding} $\beta=-0.125$, n.s.). We did find a significant increase in \textit{Logicality} ($\beta=0.216, p=0.013$), indicating that participants recognized the improved structure.

This question was guided by two insights from related literature: the multiple attention points that are native to esports (e.g., multiple battles happening simultaneously) and how Information Asymmetry can make a narrative more engaging. We attempted to address both of these through the \textit{Event Structuring} module, which attempted to discuss parallel events together, and to order events in a way that created information asymmetry (e.g., by emphasizing that a player may not be aware of a strategy). However, the results suggest that they are not sufficient to create a more engaging narrative through our approach in live text commentary.

\subsubsection{Logicality Without Engagement}
While \textit{Event Structuring} did succeed in making the text significantly more \textit{Logical} ($\beta=0.216, p=0.013$), this logicality did not translate into enjoyment.

The qualitative feedback suggests the system's approach to structuring was too shallow: it grouped related events, but did not create meaningful connections between them. A key missing link is state. Structuring a narrative requires more than reordering; it requires continuity and explicit causal connections. For example, grouping battles into a ``western front'' arc only works if the reader understands why the west matters. While the system provided some context (e.g., describing key resource locations), it often did not explain why these locations mattered in the match strategy.

Two feedback items illustrate this gap: readers were confused by missing causal explanations, and one example (Aztec Monks and otherwise-unavailable Cavalry) relied on an implicit conversion link that the text did not state. To create coherent narratives, the system needs to make these strategic connections explicit.

\subsubsection{Limits of Prompt-Based Structure}
Technically, this highlights the limitations of our prompt-based approach to structure. We relied on the LLM to infer the narrative thread from a structured JSON list of events. The results suggest this approach is too passive. A robust solution to RQ3 requires explicit narrative planning: the system must decide on a theme for a paragraph (e.g., "The collapse of the eco") and force the realization layer to use transitional phrases that highlight this theme. Mere chronological reordering is too passive to create a narrative arc.

\subsubsection{Structure Without Story}
The answer to RQ3 is somewhat inconclusive and relies in large part on the qualitative feedback. We can use high-fidelity data to structure events into logical sequences, but logicality alone does not yield a compelling narrative: true narrative coherence is a semantic property, not just chronological and salience-based clustering.

Connecting sequential event descriptions into a meaningful story requires future systems to go beyond sorting algorithms and adopt explicit narrative planning that organizes events around shared thematic or causal threads, determining the plot before sorting the scenes. This may be partly a mismatch of our chosen modality (live text commentary), but even then an automated system likely needs narrative planning supported by state tracking over longer time horizons.

\subsection{RQ4: Reporting Strategies and Viewer Content Preferences}
\textbf{RQ4: What reporting strategies are used by esports commentators and which content types do their viewers find most relevant?}

The domain analysis (Chapter~\ref{chap:prelim}) and the preliminary viewer study (Chapter~\ref{chap:viewer-study}) together answer RQ4 from two angles: what casters do, and what viewers want. There is strong agreement between the two, and both directly shaped the SAGE design.

\subsubsection{What Commentators Do}
The cross-domain analysis (RQA1--3) found that Age of Empires II (AoE2) shoutcasting closely mirrors StarCraft II (SC2) in its information-foraging structure. Casters follow the A--E--P+S loop: Actuator-level events (battles, production, movement) serve as cues; casters then enrich these with Environment (unit positions, map state) and Performance (economic comparisons) information, occasionally drawing on Sensor data (scouting). The intelligibility type distribution, dominated by \textit{What} (52\%), followed by \textit{What Could Happen} (17\%) and \textit{How To} (9\%), also closely matches SC2 \cite{dodge2018}, confirming that the core shoutcasting pattern is an RTS-genre phenomenon rather than game-specific.

At the information source level, Event Descriptions and State Updates are grounded in \textit{Descriptive Information} (direct game state), while Strategy Discussion and Prediction draw on \textit{Domain} and \textit{Analytical Knowledge}: the color-caster layer. The narrative structure is cyclical: an introduction phase, followed by a cue-driven event loop interspersed with state updates and knowledge-sharing during lulls.

\subsubsection{What Viewers Want}
The viewer study (RQB1--3) confirmed and sharpened this picture. The most valued content type was battles framed by their economic consequences (\textit{BattleEcoKills}, EMM = 4.35). This directly reflects the A--E--P+S logic: the event (Actuator) matters most when its impact on resources (Performance) is explained. Map layout introductions and state updates (age-ups, military tech) were also rated highly, while player introductions and predictions were not wanted.

Content-wise, viewers prioritized strategic understanding and skill appreciation over pure battle action. They wanted to know \textit{why} an event mattered, not just \textit{what} happened, and requested economy and win-condition context that was absent from the prototype. Coherence was achievable (participants reported good understanding), but required scouting updates and state anchors to connect events meaningfully.

\subsubsection{How This Shaped SAGE}
Both sets of findings translated directly into the SAGE architecture. The \textit{Expert Insights} module operationalizes the color-caster layer identified in the domain analysis: it adds \textit{Analytical} and \textit{Domain Knowledge} around detected events, explaining strategic significance rather than just describing game state. This directly targets the viewer preference for \textit{why} over \textit{what}. The \textit{Event Structuring} module addresses the A--E--P+S loop's concurrency challenge: parallel events are grouped and ordered to support comprehension, mirroring the caster strategy of prioritizing high-salience cues over exhaustive coverage.

The viewer study findings also shaped which hypotheses were testable. The finding that battle events with economic consequences are most valued grounded our choice of battle-detection as the primary event type. The finding that coherence requires state anchors motivated our location-naming and state-update components.

The evaluation results close the loop: \textit{Expert Insights} produced significant gains on \textit{Competitive Nature}, \textit{Skill Improvement}, and \textit{Vicarious Sensation}, precisely the skill-centric dimensions that the domain analysis and viewer study identified as the core value of analytical commentary. The struggle of \textit{Event Structuring} to improve engagement beyond logicality reflects the gap identified in the viewer study: structural ordering helps, but causal explanation (the \textit{why}) is what drives understanding and engagement.

\section{Relation to Prior Literature}

This research spans Esports Analytics, Data-to-Text Generation, and Human-Computer Interaction, particularly work relying on Large Language Models. Our findings extend established work in these areas and help clarify practical design targets.

\subsection{The Shoutcasting Model}
Existing literature on shoutcasting \cite{li2020, kempecook2019} emphasizes the dual role of play-by-play (action) and color (analysis) commentary. Our results support this split, but the contrast between \textit{Expert Insights} and the baseline suggests color commentary carries more weight than play-by-play in text-based media.
\begin{enumerate}[noitemsep,topsep=0pt]
  \item \textbf{Play-by-play}: The limited performance of our \textit{baseline} and \textit{Event Structuring} conditions challenges the utility of pure \textit{play-by-play} text generation. While existing research \cite{taniguchi2019, wang2021} focused on accurate event description without access to external information sources and our texts scored highly for structure and fluency, participant feedback labeled such descriptions as boring. The user does not need to be told ``Player A killed Unit X''; they need to be told ``Player A is winning because they killed Unit X.''
  \item \textbf{Color}: The success of \textit{Expert Insights} mirrors the Color Caster role. This aligns with \citet{sjoblom2017} and \citet{hamari2017}'s finding that learning is a primary motivation. In text, where the visual excitement is absent, the value proposition shifts toward explanation: users want to know \textit{why} events matter, not just \textit{what} happened.
\end{enumerate}

\subsection{Information Sources and Cognitive Load}
Our preliminary domain analysis (Chapter~\ref{chap:prelim}) identified a taxonomy of five information sources used by human casters: \textit{Descriptive}, \textit{Background}, \textit{Domain}, \textit{Analytical}, and \textit{Experiential}. The evaluation results can be interpreted through this taxonomy in relation to the PEAS and IFT models \cite{dodge2018}:
\begin{itemize}[noitemsep,topsep=0pt]
  \item \textbf{Descriptive}: Pure \textit{Descriptive Information} (e.g., unit counts, building completions) was not a significant driver of engagement. However, it provides a basis for low-level understanding (Logicality, Understanding) and actuator-level cues.
  \item \textbf{Analytical}: The success of \textit{Expert Insights} validates the importance of \textit{Analytical Knowledge} (derived insights) and \textit{Domain Knowledge} (meta-game context). By synthesizing multiple descriptive signals (e.g., resources + unit composition) into a single analytical statement (``Player A is rushing''), the system reduced cognitive load, effectively using the Environment and Performance layers of PEAS to communicate state.
\end{itemize}

Our text-based system does not incorporate \textit{Experiential Knowledge} (e.g., anecdotes) or rich \textit{Background Information} (e.g., player history). In human casting, these sources help fill lulls and build narrative tension. Their absence likely contributed to lower \textit{Entertaining Nature} and \textit{Dramatic Nature} scores, and may be exacerbated by the live text(-only) format.

This is consistent with \citet{zheng2025}'s proposed system, which uses a similar taxonomy to explain their design choices.

\subsection{The Limits of Automated Metrics}
Automated similarity-based metrics \cite{papineni2002, lin2004, banerjee2005} frequently used in NLG were not applicable here: they require reference texts (human-written commentary for the same events) to compare generated output against. No such corpus exists for AoE2 live commentary. This is not an unusual constraint; it is a common situation in new or specialized domains, and it is precisely the setting where human evaluation is most necessary \cite{vanderlee2019}. Our results illustrate why: the generated texts were rated as readable and coherent by participants, yet engageability remained low across most dimensions. A metric measuring surface similarity to a reference would not have captured this gap: the system's problem was not fluency but the absence of analytical depth. This is precisely why NLG evaluation should be grounded in what the target audience actually values, particularly when the goal is engagement rather than information transfer \cite{vanderlee2021}.

\subsection{Public Perception of AI}
Since many recent automated game commentators were built using LLMs \cite{wang2024,cook2024} it is important to also consider the public perception of AI. During this project a clear negative bias emerged, and this likely affected our results. What would have been novel in 2021 may be met with skepticism by 2025; a negative bias against AI-generated content may have lowered engagement scores across conditions. If distrust of AI is partly driven by uncertainty about how outputs are generated, then making the factual grounding of insights explicit (showing that claims derive from real game data rather than hallucination) may help. Our data supports this: participants who assumed the system was wrong about Aztecs fielding cavalry were in fact incorrect: the system reported a real game state (cavalry obtained via Monk conversion), but the system failed to explain the causal link, leaving the correct claim unverifiable to the reader. This suggests LLM-based systems should not be deployed as black boxes: transparency about data sources and reasoning is important for building trust in generated insights.

\section{Conclusion}
\label{sec:conclusion}

This thesis set out to investigate whether commentator-inspired engagement strategies, implemented in a data-driven pipeline, can improve the engageability of automated esports live text commentary. The work unfolded in four phases.

We began with a \textbf{domain analysis} (Chapter~\ref{chap:prelim}), adapting SC2 shoutcasting research to AoE2 to understand what professional commentators do and what information they draw on. The analysis confirmed that AoE2 shoutcasting follows the same A--E--P+S information-foraging pattern as SC2 \cite{dodge2018}: Actuator-level events trigger cue-driven commentary, enriched with Environment and Performance context. This gave us a framework for event detection and a taxonomy of information sources (Descriptive, Domain, Analytical) that could ground a prototype design. A \textbf{preliminary viewer study} (Chapter~\ref{chap:viewer-study}) then validated which of these content types actually matter to viewers. The clearest finding was that battles framed by their economic consequences, which are the most impactful in AoE2, are the most valued content, and that viewers want analytical context, \textit{why} an event matters, not play-by-play description or predictions.

These findings directly shaped the \textbf{SAGE system} (Chapter~\ref{chap:methods}). The \textit{Expert Insights} module operationalizes the color-caster side: it adds Domain and Analytical Knowledge around detected events, explaining strategic significance. The \textit{Event Structuring} module addresses the challenge of concurrent events identified in the domain analysis, grouping and ordering events to support comprehension rather than reporting them in raw chronological order. The system uses CaptureAge's high-fidelity data, going beyond the low-fidelity event logs most prior work relies on, and combines rule-based event detection with LLM-based realization.

Finally, an \textbf{evaluation} with $N=33$ AoE2 viewers (Chapter~\ref{chap:eval}) tested whether these strategies improve engagement. The results confirmed the direction of the design: \textit{Expert Insights} produced significant positive effects on skill-centric engagement dimensions (\textit{Competitive Nature}, \textit{Skill Improvement}, \textit{Vicarious Sensation}), while \textit{Event Structuring} improved perceived logicality but not engagement. The ceiling on absolute engagement scores reflects the text-only modality and survey context, not a rejection of the approach. The consistent preference for analytical insight over description, visible across both the preliminary study and the main evaluation, validates the core design decision to prioritize \textit{why} over \textit{what}.

\section{Contributions and Design Principles}

This work contributes to the fields of Human-Computer Interaction (HCI) and Natural Language Generation (NLG) through both technical artifacts and empirical findings.

\subsection{Relevance of Classical Design Frameworks}
We presented a modular framework for automated commentary based on the PEAS/IFT\cite{russell1995, pirolli1999} framework (Performance, Environment, Actuators, Sensors + Information Foraging Theory) \cite{dodge2018}. The IFT framework proved beneficial, and PEAS in the form of E-A-P+S provided the basis for the system design. It uses event detection to identify event cues, enriches these events with context and insights, and realizes them as natural language text. This architecture balances rule-based grounding with the flexibility of Large Language Models (LLMs). This core architecture, still inspired by \citet{reiter2000}: \textit{Event Detection} $\rightarrow$ \textit{Enrichment} $\rightarrow$ \textit{Narrative Realization}, remains a reusable template for turning dense temporal data into concise explanations, even as LLM capabilities grow. This design is also highly modular, allowing different modules (in our case \textit{Expert Insights} and \textit{Event Structuring}) to be injected to test specific strategies.

\subsection{Technical Considerations}
Two implementation-level constraints are likely relevant for other event-based, high-fidelity temporal systems as well. The scale and granularity of the data alongside the live operation requirement requires deliberate design considerations. First, when interpretation is inevitably delayed (e.g., due to clustering over time), systems should keep immutable, time-indexed snapshots to avoid accidental ``time travel'' and preserve what was true at the moment of the event; while perhaps obvious at first glance, this consideration should be made early in the design. Second, selective retrieval matters for context-window limits, latency, and cost alike. Blindly sending context and letting the LLM decide what is relevant is not a sustainable strategy, even during prototype phases: it leads to dropped context and high inference costs. The convenience of LLMs becomes a trap, and tuning becomes a game of whack-a-mole. Another finding worth noting is that LLMs can capture domain-specific knowledge in fuzzy form (e.g., identifying a rush strategy from multiple signals) that can be operationalized into deterministic rules, running efficiently and reliably without prompting the LLM for every inference. This reduces inference costs and latency, and produces a more explainable system with transparent, humanly verifiable inputs.

\subsection{Automated Commentary}
A clear design principle emerges from the results: \textit{Insight over Description}. Systems should prioritize analytical insights (e.g., ``Player A is rushing due to low resources'') over raw event description. We had some success by supplying LLMs with structured directives on how certain situations could warrant a specific response. The LLM was able to use this information to generate insights that were not explicitly prompted or designed. The use of supportive systems that would provide additional context to the LLM, such as the \textit{Expert Insights} module, can significantly improve the engageability of the generated text.

\subsection{Dataset}
In addition, we release two datasets (the KotD4 dataset used for this research, and a much larger variant of the dataset using modern matches and introducing an additional 1-second window resolution) to further support research in high-fidelity game replays at a fidelity and scale rare for this domain alongside the \textit{SAGE} source code (Appendix~\ref{app:corpus}, Appendix~\ref{app:source-code}).

\section{Limitations and Validity}

Several constraints frame the interpretation of these results. First, the \textit{live text commentary} format (a purely textual medium) likely imposes a ceiling on engagement scores. Because this format is novel for \textit{Age of Empires II} (AoE2), we lacked a direct baseline against human-authored live text; human commentary is primarily auditory, and written summaries are typically retrospective rather than live. In a real-world product, such text would likely complement a video player or minimap rather than stand alone. Our user study also presented the commentary in a static offline survey, which may further dampen dimensions like \textit{Dramatic Nature}. The \textit{low} scores should thus be interpreted as a reflection of the \textit{modality}, not necessarily a rejection of the \textit{quality} of the insights.

Second, while our sample size and condition count ($N=33\times4$) is comparatively large for validation studies in this domain (which often lack human evaluation entirely), it remains statistically modest and primarily powered to detect large effects. Subtle benefits of features like \textit{Event Structuring} may therefore have gone undetected. Recruitment targeted highly engaged AoE2 players, reflected in a high average skill (Elo $M=1600+$), limiting generalizability to casual viewers who might prefer simpler play-by-play narration. These factors limit granularity, but do not harm the internal validity of the primary finding: the robust directional preference for \textit{Expert Insights} is a strong signal for the design of future systems.

Third, our reliance on rule-based event detection introduced limitations in precision and recall. Heuristic rules occasionally misclassified complex interactions (e.g., merging nearby fights into a longer battle), and we do not parse all data (e.g., sensoric data such as what players know about each other's strategies). This highlights a classic trade-off: while rule-based systems offer control, capturing the fluidity of high-fidelity gameplay likely requires more adaptive, machine-learning-driven detection logic.

Our evaluation results should also be situated within rapidly shifting public opinion on AI. This project began in 2021, when large language models were still novel; by 2025, the proliferation of low-quality AI-generated content contributed to generalized skepticism toward machine-generated text. Beyond novelty, users may approach AI outputs expecting errors, \textit{hallucinations}, or a lack of originality.

\section{Future Research}
The findings suggest three follow-on research directions. Each follows directly from limitations we hit, and from what this work leaves behind: a high-fidelity dataset at unprecedented scale, a working pipeline validated with real viewers, and a community with direct access to competitive AoE2.

\subsection{The SAGE Pipeline}

The SAGE architecture represents a single vertical slice, and there are many avenues for improvement. The most direct next step is to extend the Signal Analysis and Event Detection modules to cover a broader range of event types beyond battles, and to improve the precision of existing detectors. Statistical and machine learning approaches to event detection may offer more robustness than the current rule-based heuristics. Related work that relies primarily on end-to-end generation \cite{wang2021,zhang2024} may also benefit from a pipeline approach, something proposed but not implemented by \citet{zheng2025}. Classical pipeline architectures \cite{reiter2000} remain instructive here, and exploring how they integrate with modern LLM realisers is a promising direction, including using fuzzy logic to handle the inherently gradual nature of battle onset \cite{ramos2019}.

\subsection{The SAGE Dataset}
Alongside this thesis we release the SAGE-dataset (Appendix~\ref{app:corpus}), a flexible corpus building on the KotD4 matches used during this research consisting of 173 minute-by-minute dataframes representing a complete game state. During this research we found that the 60-second snapshot resolution we used introduces attribution errors. E.g., when a unit upgrade occurs within a window, we default during realisation to the upgraded state, which may misrepresent what was visible at the time. We mitigated this by also exporting a 1-second resolution version, which has the potential to reduce attribution errors from complete misses to small timing offsets, but this dataset was not used in the current evaluation scope. Additionally, the KotD~IV set does not reliably capture explicit player commands due to a technical limitation of the data format at the time, limiting the signal available for supervised training.

A 3,100-match dataset (May 2026), also released alongside this thesis, addresses several of these constraints. It includes player commands and Fog of War, and covers a broader range of maps and civilization matchups, includes the many civilizations and mechanics added since 2021, and provides the scale needed to find patterns that a 173-match corpus cannot reliably surface.

Notably, all matches in both datasets are replayable in CaptureAge, meaning the data can be cross-referenced against visual game state at any timestamp. To our knowledge, no comparable publicly available \cite{synnaeve2017} dataset of this resolution \cite{aoe2ocel2024} and coverage exists for similar games. In particular a dataset like this can realistically support the statistical and machine learning approaches we mentioned in the previous sub-section.

\subsection{The Knowledge-Latency Trade-off}

The Aztec cavalry example from the evaluation illustrates a fundamental limitation that goes beyond event detection: the realization stage lacked access to the causal knowledge needed to explain what it was reporting. The system correctly identified that an Aztec player had cavalry (a rare and unusual state) but did not explain that Aztecs cannot train cavalry natively, and that the only way to obtain them is through Monk conversion. Without that chain of reasoning (\textit{Aztec} $\rightarrow$ \textit{no cavalry training} $\rightarrow$ \textit{must have converted via Monks}), a reader with domain knowledge correctly identifies the statement as suspicious, even though it was factually accurate. While it is tempting to blame the reader, it is the system's responsibility to make the reasoning process transparent and verifiable, even if that simply means nudging the reader's attention to the relevant data (e.g., ``Player A has Aztecs, and has converted 3 Cavalry with Monks'').

Addressing this requires a knowledge base of domain-specific facts that can be retrieved at realization time. A promising direction is to pre-generate these facts using an LLM (civilization bonuses, unit counters, strategic implications of specific matchups), review them for correctness, and store them in a structured database for lookup. We already observed that LLMs can produce such representations for this domain (e.g., identifying build-orders and encoding them as fuzzy rules): these were used directly in the Expert Insights module after manual review. Extending this to a broader, indexed knowledge base that is queried at generation time would align well with Retrieval-Augmented Generation (RAG) approaches \cite{lewis2020}, and is consistent with prior work that uses structured knowledge to support generation, whether by retrieving contextual story material in sports commentary \cite{lee2013} or by grounding template-based reports in an external fact database at generation time \cite{gatti2018}.

Live commentary imposes real-time constraints: a multi-step agentic pipeline that retrieves facts, chains inferences, and realizes text is not feasible if it adds several seconds of delay per update. In practice, the system can distinguish between \textit{pre-generated} knowledge (static facts that can be retrieved quickly) and \textit{dynamic inference} (contextual reasoning triggered by a specific game state). The fuzzy logic rules generated as a side effect of building SAGE suggest that pre-generation is more tractable than it appears: much of the domain knowledge an LLM would reason about at generation time can instead be materialized into lookup structures in advance, trading reasoning time for retrieval time.

\section{Final Remarks}
At the core of this project we wanted to create a system that could generate engaging esports commentary. This requires a human-centered approach to design and evaluation; however, most research in this field is centered around end-to-end generation where the groundwork of understanding the domain and the audience is often skipped, at both design and evaluation stages.

Before a single module of SAGE was designed, we studied what AoE2 shoutcasters actually do: how they move between event cues and analytical context, which information sources they draw on, and how their structure mirrors the SC2 pattern established by prior work. Before any hypothesis was tested, we asked viewers what they valued: which event types mattered, how to balance action against strategy, and whether event-driven summaries could form coherent narratives. Without them, SAGE would have had no principled answer to what \textit{Expert Insights} to add, which events deserved priority, or what a meaningful commentary text should look like.

The human evaluation closed this loop. The significant gains on skill-centric engagement dimensions were not a lucky result; they followed directly from what the domain analysis and viewer study had told us viewers valued. \textit{Event Structuring}'s inability to improve engagement beyond logicality was clear in hindsight: the viewer study had already shown that structural ordering helps, but adding (or showing) explanation is what drives understanding, something they clearly stated, and the decision to add structuring was not based on the audience's direct feedback, but rather different angles from other research indicating structure can create engagement.

We did attempt this route: aligning a T90Official shoutcast directly with the CaptureAge data frames, using the camera at each timestamp to infer which game state elements were visible when specific words were spoken. This would, in principle, allow training a data-to-text model on (observed state, spoken commentary) pairs without manual annotation. It is, however, a project of a different scale and scope, not realistically achievable for smaller projects in niche domains. And even then, there is no guarantee that viewers would appreciate the content produced without being part of the validation process.

This is a broader point for data-to-text systems in specialized domains. When no annotated corpus exists and no direct precedent applies, studying the experts who have internalized the domain's narrative conventions, and validating with the audience who defines what the output should be, provides a foundation to build upon. The most transferable contribution may be the approach itself: ground the design in human behavior before building the system, and evaluate against human expectations before drawing conclusions. In a domain as specialized and underexplored as esports commentary, this is what made the work possible. This is not new, but the push toward LLM-based systems has led to a resurgence of end-to-end approaches that skip these steps. We think this work illustrates how taking the human into account end-to-end, rather than the data, is what allows us to build systems that are technically functional and meaningfully engaging to real users.
